\SourceLeaf{316}%
Die hiemit wiederholte Rechnung ergibt
\[
\begin{aligned}
\beta&=3407''\,9\\
B&=51^0\,17'\,39''\,7\\
\tau&=420''\,55\\
T&=5^0\,38'\,51''\,5
\end{aligned}
\]
Mit diesem Werthe von \(T\) wird nun die schärfere Rechnung angefangen und dabei zugleich die logarithmische Correction mit zugezogen. Es findet sich. in Einheiten der siebenten Decimale,
\[
\begin{aligned}
\mu rr&=99.76\\
\mu\lambda\lambda&=2.47\\
\mu\tau\tau&=1.50
\end{aligned}
\]
folglich
\[
\begin{aligned}
\mu\lambda\lambda+\tfrac12\mu\tau\tau&=+3\\
\mu rr+\tfrac12\mu\tau\tau&=+101\\
-\tfrac12\mu rr+\tfrac12\mu\lambda\lambda&=-49
\end{aligned}
\]
und
\[
\begin{aligned}
\log\beta&=3.5324974\\
\log b&=3.5324977\\
b&=3407''\,9852\\
B&=51^0\,17'\,39''\,6419\\
\log\tau&=2.6238492\\
\log t&=2.6238593\\
t&=420''\,5904\\
T&=5^0\,38'\,51''\,4752
\end{aligned}
\]
Eine nochmalige Wiederholung der Rechnung mit diesem Werthe von \(T\) bringt bei \(b\) gar keine Änderung hervor, und \(t\) verwandelt sich in \(420''\,5898\). Man erhält daher
\begin{center}Breite des Punkts Inselsberg\end{center}
\[
B^0-\beta=50^0\,49'\,15''\,6493
\]
\begin{center}Azimuth der Dreiecksseite Inselsberg-Brocken\end{center}
\[
T^0-\tau+180^0=185^0\,35'\,21''\,1806
\]
